% --------------------------------------------------------------------------
%  Common Styles and Formattings for Kharazmi University Proposal
%  تنظیمات و بسته‌های مورد نیاز برای پروپوزال دانشگاه خوارزمی
%  دانشکده فنی و مهندسی
% --------------------------------------------------------------------------

\usepackage[table,xcdraw]{xcolor}
\usepackage{amssymb,amsmath}
\usepackage{graphicx}
\usepackage{geometry}
\usepackage{fancyhdr}
\usepackage{multicol}
\usepackage{enumitem}
\usepackage{multirow}
\usepackage{array}
\usepackage{tabularx}
\usepackage{booktabs}
\usepackage{makecell}
\usepackage{colortbl}
\usepackage{float}
\usepackage{caption}
\usepackage[colorlinks,linkcolor=blue,citecolor=blue,urlcolor=blue]{hyperref}

% -------------------- Page Layout (ابعاد صفحه مطابق قالب رسمی خوارزمی) --------------------
\geometry{
    a4paper,
    top=2.0cm,
    bottom=2.2cm,
    right=2.5cm,
    left=2.2cm,
    headheight=0pt,
    headsep=0pt,
    footskip=22pt
}

% فاصله‌گذاری خطوط و پاراگراف‌ها
\linespread{1.3}
\setlength{\parskip}{0.35em}
\setlength{\parindent}{0.8cm} % فرورفتگی ابتدای پاراگراف‌ها

% -------------------- Headers & Footers (بدون سربرگ، شماره صفحه پایین-چپ) --------------------
\pagestyle{fancy}
\fancyhf{}
\renewcommand{\headrulewidth}{0pt}
\renewcommand{\footrulewidth}{0pt}
\fancyfoot[L]{\thepage} % شماره صفحه در پایین سمت چپ با اعداد فارسی

\fancypagestyle{plain}{
    \fancyhf{}
    \renewcommand{\headrulewidth}{0pt}
    \renewcommand{\footrulewidth}{0pt}
    \fancyfoot[L]{\thepage}
}

% -------------------- تعاریف ستون‌های جداول (تراز وسط عمودی و افقی) --------------------
% در زی‌پرشین، raggedright متن را راست‌چین و raggedleft متن را چپ‌چین می‌کند
\newcolumntype{C}[1]{>{\centering\arraybackslash}m{#1}}
\newcolumntype{R}[1]{>{\raggedright\arraybackslash}m{#1}}
\newcolumntype{L}[1]{>{\raggedleft\arraybackslash}m{#1}}

% -------------------- نمادهای بالت در فهرست‌ها (جلوگیری از کاراکترهای نامعتبر در فونت) --------------------
\setlist[itemize,1]{label=\raisebox{0.15ex}{\small$\bullet$}}
\setlist[itemize,2]{label=\raisebox{0.15ex}{\footnotesize$\circ$}}
\setlist[itemize,3]{label=\raisebox{0.15ex}{\tiny$\bullet$}}
\renewcommand{\labelitemi}{\raisebox{0.15ex}{\small$\bullet$}}
\renewcommand{\labelitemii}{\raisebox{0.15ex}{\footnotesize$\circ$}}

% -------------------- تیترهای بخش‌ها (Kharazmi Headings) --------------------
% تیترها دقیقاً سایز ۱۲ بولد با خط زیرِ خود متن، بدون شماره و تراز راست
\newcommand{\khusection}[1]{%
    \par\vspace{1.5em}%
    \noindent{\normalsize\bfseries\underline{#1}}%
    \par\vspace{0.6em}%
}

% -------------------- Checkbox Symbols --------------------
\newcommand{\uncheckedbox}{\raisebox{-0.1ex}{\large$\square$}}
\newcommand{\checkedbox}{\raisebox{-0.1ex}{\large$\boxtimes$}}

\newcommand{\boxAdadi}{\ifGrantAdadi\checkedbox\else\uncheckedbox\fi}
\newcommand{\boxAzmayeshgahi}{\ifGrantAzmayeshgahi\checkedbox\else\uncheckedbox\fi}
\newcommand{\boxRouzaneh}{\ifGrantRouzaneh\checkedbox\else\uncheckedbox\fi}
\newcommand{\boxPardis}{\ifGrantPardis\checkedbox\else\uncheckedbox\fi}

% -------------------- Persian & Algorithms --------------------
\usepackage{algorithmicx,algorithm}
\usepackage[noend]{algpseudocode}

\usepackage{xepersian}
%------------------------ Algorithm ------------------------------------

\newenvironment{الگوریتم}[1]
	{\bigskip\bigskip\begin{algorithm}\caption{#1} \label{الگوریتم: #1}\vspace{0.5em}\begin{algorithmic}[1]}
	{\end{algorithmic}\vspace{0.5em}\end{algorithm}\bigskip}
	

\renewcommand{\algorithmicfor}{{به ازای}}
\renewcommand{\algorithmicwhile}{{تا وقتی}}
\renewcommand{\algorithmicdo}{\hspace{-.2em}:}
\renewcommand{\algorithmicif}{{اگر}}
\renewcommand{\algorithmicthen}{\hspace{-.2em}:}
\renewcommand{\algorithmicelse}{{در غیر این صورت:}}
%\renewcommand{\algorithmicelsif}{{در غیر این صورت اگر: }}
\renewcommand{\algorithmicreturn}{{برگردان}}
\renewcommand{\algorithmiccomment}[1]{$\triangleleft$ \emph{#1}}
\renewcommand{\algorithmicrequire}{\textbf{ورودی:}}
\renewcommand{\algorithmicensure}{\textbf{خروجی:}}

\newcommand{\اگر}{\If}
\newcommand{\وگرنه}{\Else}
\newcommand{\وگر}{\ElsIf}
\newcommand{\پایاناگر}{\EndIf}
\newcommand{\بهازای}{\For}
\newcommand{\پایانبهازای}{\EndFor}
\newcommand{\تاوقتی}{\While}
\newcommand{\پایانتاوقتی}{\EndWhile}
\newcommand{\دستور}{\State}
\newcommand{\دستورک}{\Statex}
\newcommand{\توضیحات}{\Comment}
\newcommand{\برگردان}{\Return}
\providecommand{\ورودی}{\Require}
\providecommand{\خروجی}{\Ensure}


% -------------------- Fonts --------------------
\settextfont[
    Scale=1.0,
    Extension=.ttf, 
    Path=styles/fonts/,
    BoldFont=BNaznnBd,
    ItalicFont=XB NiloofarIt,
    BoldItalicFont=XB NiloofarBdIt,
    Script=Persian
]{BNazanin}

\setlatintextfont[
    Scale=1.0,
    Extension=.ttf, 
    Path=styles/fonts/,
    BoldFont=timesbd,
    ItalicFont=timesi,
    BoldItalicFont=timesbi
]{times}

\setmathdigitfont[
    Scale=1.0,
    Extension=.ttf, 
    Path=styles/fonts/,
    BoldFont=XB Yas Bd,
    ItalicFont=XB Yas It,
    BoldItalicFont=XB Yas BdIt
]{XB Yas}

\defpersianfont\titrfont[
    Scale=1.0,
    Path=styles/fonts/,
    Extension=.ttf,
    BoldFont=BNaznnBd
]{BNaznnBd}

% -------------------- تعاریف ستون‌های انگلیسی و ماکروهای کمکی دو زبانه --------------------
% ستون استاندارد برای متن‌های انگلیسی چپ‌چین در جداول با فونت لاتین و جهت چپ به راست
\newcolumntype{E}[1]{>{\begin{latin}\setlength{\parskip}{0pt}\raggedright\arraybackslash}m{#1}<{\strut\par\end{latin}}}

% -------------------- حذف تیتر خودکار و تکراری «مراجع» از thebibliography --------------------
\makeatletter
\AtBeginDocument{%
  \renewenvironment{thebibliography}[1]
       {\list{\@biblabel{\@arabic\c@enumiv}}%
             {\settowidth\labelwidth{\@biblabel{#1}}%
              \leftmargin\labelwidth
              \advance\leftmargin\labelsep
              \@openbib@code
              \usecounter{enumiv}%
              \let\p@enumiv\@empty
              \renewcommand\theenumiv{\@arabic\c@enumiv}}%
        \sloppy
        \clubpenalty4000
        \@clubpenalty \clubpenalty
        \widowpenalty4000%
        \sfcode`\.\@m}
       {\def\@noitemerr
         {\@latex@warning{Empty `thebibliography' environment}}%
        \endlist}%
}
\makeatother

% -------------------- Footnotes & Bilingual Helpers --------------------
\autofootnoterule
\let\origfootnoterule\footnoterule
\renewcommand{\footnoterule}{\vfill\origfootnoterule}
\renewcommand{\thefootnote}{\arabic{footnote}}
\newcommand{\پاورقی}[1]{\LTRfootnote{#1}}
\newcommand{\کد}[1]{\lr{\ttfamily #1}}
\newcommand{\انگلیسی}[1]{\lr{#1}}
\newcommand{\فارسی}[1]{\rl{#1}}

